%----------------------------------------------------------------------
%  Chapter 86 — Boundary Condition Imposition and Incremental Control
%                Architecture for Multiscale Analysis
%----------------------------------------------------------------------
\chapter{Boundary Condition Imposition and Incremental Control
         Architecture for Multiscale Analysis}
\label{ch:bc-incremental-multiscale}

This chapter provides a detailed exposition of how Dirichlet boundary
conditions are imposed, how displacement (and load) increments are
controlled through the incremental stepping engine, and how these
mechanisms interact with the FE\textsuperscript{2} multiscale coupling
framework.  It also identifies a fundamental gap in the current
architecture: the absence of a complete \emph{physical state transfer}
protocol between macro- and micro-scale models, and proposes extensions
that would inject reconstructed stress tensors, conjugate kinematic
measures, and material history variables into the sub-model, thereby
improving both the fidelity of the physics and the convergence of the
nonlinear sub-model solver.


%======================================================================
\section{Dirichlet BC Imposition via PETSc Section Elimination}
\label{sec:dirichlet-imposition}
%======================================================================

Unlike many finite element codes that employ penalty methods or
Lagrange multipliers, \texttt{fall\_n} uses the native PETSc DMPlex
\emph{section constraint} mechanism, which implements algebraic
elimination of constrained degrees of freedom.  The approach proceeds
in three phases.

% ─────────────────────────────────────────────────────────────────────
\subsection{Phase 1: Constraint Declaration}
\label{ssec:bc-declaration}

Before calling \verb|setup()|, the user (or the FE\textsuperscript{2}
framework) declares constraints through the \texttt{Model} API:

\begin{lstlisting}
// Fixed base (all 6 DOFs = 0)
model.constrain_node(0, {0, 0, 0, 0, 0, 0});

// Pre-declare controlled DOF (value will change during analysis)
model.constrain_dof(top_node, 0, 0.0);

// Selector-based: roller support on x=0 plane (only u_x fixed)
model.constrain_dof_where(PlaneSelector<3>{0, 0.0}, 0, 0.0);
\end{lstlisting}

These methods store constraints in the internal map
\begin{equation}
  \texttt{constraints\_} :
    \underbrace{\mathit{plex\_id}}_{\text{DMPlex point}}
    \;\longrightarrow\;
    \bigl(\,
      \underbrace{\mathbf{i}}_{\substack{\text{sorted DOF}\\\text{indices}}},\;
      \underbrace{\mathbf{v}}_{\text{prescribed values}}
    \,\bigr),
  \label{eq:constraint-map}
\end{equation}
where the DOF indices within each node are maintained in sorted order
(a requirement of \texttt{PetscSectionSetConstraintIndices}).

% ─────────────────────────────────────────────────────────────────────
\subsection{Phase 2: PETSc Section Assembly}
\label{ssec:bc-section-assembly}

During \verb|setup()|, the constraint information is written into the
PETSc \texttt{PetscSection} that governs the DOF layout:

\begin{enumerate}
  \item \textbf{Mark constrained DOF counts:}
    \begin{lstlisting}[basicstyle=\small\ttfamily]
for (auto& [plex_id, idx] : constraints_)
    PetscSectionAddConstraintDof(dof_section_, plex_id,
                                 idx.first.size());
PetscSectionSetUp(dof_section_);
    \end{lstlisting}
  \item \textbf{Set constrained DOF indices:}
    \begin{lstlisting}[basicstyle=\small\ttfamily]
for (const auto& [plex_id, idx] : constraints_)
    PetscSectionSetConstraintIndices(dof_section_, plex_id,
                                     idx.first.data());
    \end{lstlisting}
  \item \textbf{Attach section to DM and create vectors:}
    \begin{lstlisting}[basicstyle=\small\ttfamily]
DMSetLocalSection(dm, dof_section_);
DMSetUp(dm);
DMCreateLocalVector(dm, &global_imposed_solution_);
    \end{lstlisting}
\end{enumerate}

The key consequence is that the PETSc \emph{global} vectors and
matrices created from this DM automatically exclude constrained DOFs.
The global solution vector~$\mathbf{U}$ contains only the
\emph{free} DOFs~$\mathbf{u}_f$, while the prescribed values~$\mathbf{u}_c$
reside in the \emph{local} vector~\verb|global_imposed_solution_|.

% ─────────────────────────────────────────────────────────────────────
\subsection{Phase 3: Residual Assembly and Imposed Values}
\label{ssec:bc-residual}

Within each Newton iteration, the SNES residual callback reconstructs
the full displacement field by combining free and constrained DOFs:

\begin{equation}
  \mathbf{u}_{\text{total}} =
    \underbrace{\texttt{DMGlobalToLocal}(\mathbf{U})}_{\text{expand free DOFs}}
    \;+\;
    \underbrace{\mathbf{u}_{\text{imposed}}}_{\text{prescribed values}}.
  \label{eq:u-total}
\end{equation}

In code:
\begin{lstlisting}
DMGlobalToLocal(dm, u_global, INSERT_VALUES, u_local);
VecAXPY(u_local, 1.0, model->imposed_solution());  // u += u_c
\end{lstlisting}

Each element then computes its internal force vector
$\mathbf{f}_e^{\text{int}}(\mathbf{u}_{\text{total}})$ using the
complete field.  This naturally incorporates the coupling between free
and constrained DOFs
\begin{equation}
  \mathbf{R}(\mathbf{u}_f)
    = \mathbf{f}^{\text{int}}(\mathbf{u}_f + \mathbf{u}_c) - \mathbf{f}^{\text{ext}}
    = \mathbf{0},
  \label{eq:residual-with-bc}
\end{equation}
so that the $\mathbf{K}_{fc}\,\mathbf{u}_c$ contribution appears
automatically in the residual without explicit condensation.

% ─────────────────────────────────────────────────────────────────────
\subsection{Dynamic Update of Imposed Values}
\label{ssec:bc-update-imposed}

During incremental analysis, the method
\begin{lstlisting}
model->update_imposed_value(node_idx, local_dof, value);
\end{lstlisting}
performs two atomic operations:
\begin{enumerate}
  \item updates the value in the in-memory \texttt{constraints\_} map, and
  \item writes the value into
    \texttt{global\_imposed\_solution\_} via
    \texttt{VecSetValueLocal(\ldots, INSERT\_VALUES)}.
\end{enumerate}

This is $O(\log n)$ for the map update and $O(1)$ for the PETSc vector
insertion.  It is safe to call inside the incremental stepping loop
because it never mutates the DM section---only the numerical value of the
imposed displacement changes, not the DOF layout.


%======================================================================
\section{Incremental Control Policy Architecture}
\label{sec:incremental-control}
%======================================================================

The stepping engine of \texttt{NonlinearAnalysis} is parameterised by a
\emph{control policy} that determines the complete system state for a
given scalar parameter~$p \in [0,1]$.  The design is formalised as a
C++20 concept:

\begin{lstlisting}
template <typename S, typename ModelT>
concept IncrementalControlPolicy =
    requires(S& s, double p, Vec f_full, Vec f_ext, ModelT* m) {
        s.apply(p, f_full, f_ext, m);
    };
\end{lstlisting}

% ─────────────────────────────────────────────────────────────────────
\subsection{The Absolute Idempotency Invariant}
\label{ssec:idempotency}

The central design invariant is:

\begin{quote}
\itshape
For any value of $p$, calling \texttt{scheme.apply(p, \ldots)} must
fully determine the system state (external forces, imposed
displacements, and any other boundary conditions).
The call is \emph{absolute} and \emph{idempotent}: the result depends
only on~$p$, not on the history of previous \texttt{apply()} calls.
\end{quote}

This invariant makes the bisection engine trivially correct.  When a step
$[p_{\text{prev}} \to p_{\text{target}}]$ diverges, the engine:
\begin{enumerate}
  \item reverts $\mathbf{U}$ to its snapshot,
  \item reverts material internal variables to their committed state,
  \item calls \texttt{scheme.apply($p_{\text{prev}}$, \ldots)} to
    restore the boundary/load state (no incremental state to unwind),
  \item bisects:
    $p_{\text{mid}} = \tfrac{1}{2}(p_{\text{prev}} + p_{\text{target}})$,
  \item recursively solves
    $[p_{\text{prev}} \to p_{\text{mid}}]$ then
    $[p_{\text{mid}} \to p_{\text{target}}]$,
  \item up to a configurable maximum bisection depth.
\end{enumerate}

This is fundamentally different from \emph{incremental} schemes
(e.g.~OpenSees \texttt{integrator DisplacementControl}) where state
must be explicitly saved and restored during bisection.

% ─────────────────────────────────────────────────────────────────────
\subsection{Concrete Control Policies}
\label{ssec:control-policies}

Four concrete policies satisfy the concept:

\begin{table}[htbp]
\centering
\caption{Available incremental control policies.}
\label{tab:control-policies}
\begin{tabular}{@{}lll@{}}
\toprule
\textbf{Policy} & \textbf{System State at~$p$} & \textbf{Usage} \\
\midrule
\texttt{LoadControl}
  & $\mathbf{f}_{\text{ext}} = p\,\mathbf{f}_{\text{full}}$
  & Proportional loading \\
\texttt{DisplacementControl}
  & $u_i^{\text{imp}} = p\,u_i^{\text{target}}$;\;
    $\mathbf{f}_{\text{ext}} = p\,s\,\mathbf{f}_{\text{full}}$
  & Monotonic pushover \\
\texttt{CustomControl<F>}
  & User-supplied lambda
  & Cyclic protocols \\
\texttt{ArcLengthControl}
  & Crisfield cylindrical/spherical
  & Snap-through, post-peak \\
\bottomrule
\end{tabular}
\end{table}

% ─────────────────────────────────────────────────────────────────────
\subsection{Comparison with OpenSees TimeSeries}
\label{ssec:timeseries-comparison}

In OpenSees, loading is decomposed as
$\mathbf{f}(t) = \lambda(t) \cdot \mathbf{f}_{\text{pattern}}$,
where $\lambda(t)$ is a \texttt{TimeSeries} object that maps pseudo-time
to a scalar load factor, and the load pattern associates forces with
nodes.  The two abstractions are composed multiplicatively.

In \texttt{fall\_n}, the \texttt{IncrementalControlPolicy} serves a
similar role but with three key differences:

\begin{enumerate}
  \item \textbf{Scope:} An OpenSees \texttt{TimeSeries} can only scale
    loads.  A control policy can modify \emph{any} aspect of the system
    state: external forces, imposed displacements, and even model
    parameters---all within a single \texttt{apply()} call.
  \item \textbf{Idempotency:} OpenSees \texttt{TimeSeries} objects
    can be stateful (e.g.\ \texttt{PathTimeSeries} with interpolation
    over time history).  Control policies are \emph{required} to be
    idempotent, ensuring bisection correctness.
  \item \textbf{Composition:} OpenSees allows multiple
    \texttt{LoadPattern} objects active simultaneously.
    In \texttt{fall\_n}, all loading logic is composed inside a single
    lambda passed to \texttt{make\_control()}:
    \begin{lstlisting}
auto scheme = make_control([&](double p, Vec f_full,
                               Vec f_ext, auto* m) {
    VecCopy(f_full, f_ext);
    VecScale(f_ext, gravity_factor(p));
    VecAXPY(f_ext, lateral_factor(p), f_lateral);
    m->update_imposed_value(node, 0,
                            cyclic_displacement(p, amps));
});
    \end{lstlisting}
\end{enumerate}

A separate \texttt{TimeSeries} abstraction is therefore \emph{not
warranted}: the current design is both more general and provably
compatible with the bisection engine.  What \emph{is} valuable is a
library of reusable $p \to \text{value}$ functions (such as the existing
\texttt{cyclic\_displacement()} protocol), which serve the role of
composable building blocks without introducing stateful objects.


%======================================================================
\section{Cyclic Displacement Protocol}
\label{sec:cyclic-protocol}
%======================================================================

The function \texttt{cyclic\_displacement(p, amplitudes)} maps the
scalar parameter~$p \in [0,1]$ to a physical displacement (in metres)
according to a symmetric triangular wave protocol with monotonically
increasing amplitude levels:

\begin{equation}
  d(p) = \begin{cases}
    f \cdot A_k           & \text{phase 0: } 0 \to +A_k \\
    A_k\,(1 - 2f)         & \text{phase 1: } +A_k \to -A_k \\
    -A_k\,(1 - f)          & \text{phase 2: } -A_k \to 0
  \end{cases}
  \label{eq:cyclic-displacement}
\end{equation}

where $k = \lfloor s / 3 \rfloor$ is the amplitude level index,
$s = \lfloor 3\,N_{\text{levels}} \cdot p \rfloor$ is the segment
index, $f = 3\,N_{\text{levels}} \cdot p - s$ is the intra-segment
fraction, and $A_k$ is the $k$-th amplitude from the user-specified
sequence.  This function is \emph{pure} (no side effects, no internal
state), preserving the idempotency invariant of the control policy.

Three preset amplitude sequences are provided:

\begin{table}[htbp]
\centering
\caption{Preset cyclic amplitude sequences.}
\label{tab:cyclic-amplitudes}
\begin{tabular}{@{}lcl@{}}
\toprule
\textbf{Preset} & \textbf{Levels} & \textbf{Amplitudes (mm)} \\
\midrule
Legacy (20\,mm)     & 4 & $\pm 2.5,\; \pm 5,\; \pm 10,\; \pm 20$ \\
Extended (50\,mm)   & 6 & $\pm 2.5,\; \pm 5,\; \pm 10,\; \pm 20,\; \pm 35,\; \pm 50$ \\
Extended (100\,mm)  & 9 & $\pm 2.5,\; \ldots,\; \pm 65,\; \pm 80,\; \pm 100$ \\
\bottomrule
\end{tabular}
\end{table}


%======================================================================
\section{Worked Example: Cyclic Beam Validation}
\label{sec:worked-example-cyclic}
%======================================================================

The following annotated code fragment shows the complete boundary
condition setup and incremental solve for the cyclic validation
of a beam column element:

\begin{lstlisting}
// -- Phase 1: Declare constraints (before setup) ----------------
model.constrain_node(0, {0,0,0,0,0,0});      // fixed base
model.constrain_dof(top_node, 0, 0.0);       // pre-declare DOF 0
model.setup();                                // DM section locked

// -- Phase 2: Define control scheme ------------------------------
auto scheme = make_control(
    [top_node, &cfg](double p, Vec, Vec f_ext, BeamModel* m) {
        VecSet(f_ext, 0.0);                   // no external forces
        m->update_imposed_value(top_node, 0,  // DOF 0 at top
                                cfg.displacement(p));  // cyclic!
    });

// -- Phase 3: Incremental solve with bisection ------------------
bool ok = nl.solve_incremental(
    cfg.total_steps(),         // e.g. 180 steps
    cfg.max_bisections,        // e.g. 6 levels
    scheme);                   // injected policy
\end{lstlisting}

At each of the 180 uniform steps, the parameter advances as
$p_k = k/180$, and the control policy maps $p_k$ through
\texttt{cyclic\_displacement()} to obtain the physical displacement.
If SNES diverges at a given $p_k$, the bisection engine
halves the step and retries, recursively, up to the specified depth.
The scheme never needs to be ``restored'' because it is an absolute
function of~$p$.


%======================================================================
\section{Initial Physical State Imposition for FE\textsuperscript{2}}
\label{sec:initial-state-fe2}
%======================================================================

The preceding sections describe how \emph{kinematic} boundary conditions
(displacement Dirichlet) are imposed on the sub-model in each FE%
\textsuperscript{2} step.  This section examines the current state
transfer protocol and identifies critical gaps that must be addressed
to achieve a physically consistent and computationally efficient
multiscale coupling.

% ─────────────────────────────────────────────────────────────────────
\subsection{Current State: Displacement-Only Downscaling}
\label{ssec:displacement-only}

In the present implementation
(\texttt{NonlinearSubModelEvolver}, Chapter~\ref{ch:nl-sub-model-evolver}),
the global-to-local transfer consists exclusively of
\emph{displacement boundary conditions} derived from the beam's
Timoshenko kinematic reconstruction:

\begin{equation}
  \mathbf{u}^{\text{RVE}}(\mathbf{x}_{\text{face}})
    = \mathbf{u}_c + \boldsymbol{\theta} \times
      (\mathbf{x}_{\text{face}} - \mathbf{x}_c),
  \label{eq:kinematic-bc}
\end{equation}

where $\mathbf{u}_c$ is the centroidal translation,
$\boldsymbol{\theta}$ the rotation vector, and
$\mathbf{x}_c$~the section centroid, all in the beam's local frame.
These displacements are imposed on the two end faces (min-$z$
and max-$z$) of the prismatic sub-model domain.

The state transfer flow is:
\begin{enumerate}
  \item Extract \texttt{SectionKinematics} from the beam element
    at parametric coordinates $\xi = \pm 1$.
  \item Evaluate Eq.~\eqref{eq:kinematic-bc} at each boundary node.
  \item Call \verb|model.constrain_node(nid, {u_x, u_y, u_z})| for
    each boundary node.
  \item Solve the sub-model via incremental displacement control
    with 20 steps and 8 bisection levels.
\end{enumerate}

The \texttt{NonlinearSubModelEvolver} is \emph{persistent}: the
displacement vector~$\mathbf{U}$ is \emph{not} reset between
global steps, so the converged sub-model state from step~$n$ serves
as the initial guess for step~$n+1$.  Material internal variables
(crack openings, plastic strains, Menegotto--Pinto reversal points)
accumulate across the loading history.

% ─────────────────────────────────────────────────────────────────────
\subsection{The ModelState Snapshot}
\label{ssec:model-state}

The \texttt{ModelState<dim>} structure provides a PETSc-independent
snapshot of a model's solved state:

\begin{equation}
  \mathcal{S} = \Bigl\{
    \bigl\{n_i,\; \mathbf{u}_i,\; \dot{\mathbf{u}}_i\bigr\}_{i=1}^{N_{\text{nodes}}},
    \;\;
    \bigl\{e_j,\;
      \bigl\{\boldsymbol{\sigma}_{jk},\;
             \boldsymbol{\varepsilon}_{jk}
      \bigr\}_{k=1}^{N_{\text{GP},j}}
    \bigr\}_{j=1}^{N_{\text{elem}}}
  \Bigr\}.
  \label{eq:model-state}
\end{equation}

This structure captures nodal displacements and velocities, as well as
Gauss-point stress and strain tensors in Voigt notation.
The \texttt{capture\_state()} method of \texttt{Model<>} populates
this structure from the PETSc vectors and element material points,
while \texttt{apply\_initial\_displacement()} writes a captured
displacement field back as the initial imposed solution.

\paragraph{Limitation.}  The current \texttt{ModelState} stores
stress and strain but has no mechanism to inject them \emph{into} a
model's material points as initial conditions.  It captures the state
for post-processing and continuation (displacement seeding), but does
not restore the \emph{constitutive} state.

% ─────────────────────────────────────────────────────────────────────
\subsection{What Is Missing: Stress and Strain Injection}
\label{ssec:stress-injection}

In a rigorous FE\textsuperscript{2} downscaling, the macro-to-micro
transfer should include not only the kinematic field but also the
\emph{reconstructed stress tensor}~$\boldsymbol{\sigma}^{\text{macro}}$
and the conjugate kinematic measure~$\boldsymbol{\varepsilon}^{\text{macro}}$,
evaluated at each Gauss point of the sub-model from the beam's generalised
strains.

The Timoshenko kinematics provide the generalised strain measures at
the section level:
\begin{equation}
  \mathbf{e} = \bigl\{
    \varepsilon_0,\; \kappa_y,\; \kappa_z,\;
    \gamma_y,\; \gamma_z,\; \theta'
  \bigr\}^T,
  \label{eq:gen-strains}
\end{equation}
from which the 3D strain field at any point $(y,z)$ in the section can
be reconstructed:
\begin{equation}
  \boldsymbol{\varepsilon}^{\text{macro}}(y,z) =
    \varepsilon_0\,\mathbf{e}_z \otimes \mathbf{e}_z
    - \kappa_y\,z\,\mathbf{e}_z \otimes \mathbf{e}_z
    + \kappa_z\,y\,\mathbf{e}_z \otimes \mathbf{e}_z
    + \tfrac{1}{2}\gamma_y \bigl(
        \mathbf{e}_y \otimes \mathbf{e}_z
      + \mathbf{e}_z \otimes \mathbf{e}_y\bigr)
    + \tfrac{1}{2}\gamma_z \bigl(
        \mathbf{e}_x \otimes \mathbf{e}_z
      + \mathbf{e}_z \otimes \mathbf{e}_x\bigr),
  \label{eq:strain-reconstruction}
\end{equation}
and the corresponding stress tensor via the section's effective tangent:
\begin{equation}
  \boldsymbol{\sigma}^{\text{macro}}(y,z)
    = \mathbb{D}_{\text{sec}}^{-1}(\mathbf{e}) \cdot
      \boldsymbol{\varepsilon}^{\text{macro}}(y,z).
  \label{eq:stress-reconstruction}
\end{equation}

If these tensors were injected into each Gauss point of the sub-model
\emph{before} the nonlinear solve, the benefits would be:

\begin{enumerate}
  \item \textbf{Better initial guess:}  The SNES solver would start
    from a state that already satisfies equilibrium at the macro-scale,
    rather than from a zero-stress (or purely kinematic) initial
    condition.  This can dramatically reduce the number of Newton
    iterations required.

  \item \textbf{Physical consistency:}  The sub-model's material
    points would begin from a stress state consistent with the
    global solution, preventing spurious elastic unloading in the
    first Newton iteration that arises when the displacement field
    is imposed but the stress field is zero.

  \item \textbf{Compatibility with path-dependent materials:}
    Damage and plasticity models depend on the stress path, not just
    the final strain.  Injecting the macro-scale stress provides
    a better ``starting ramp'' for the local constitutive integration.
\end{enumerate}


%======================================================================
\section{Material History Mapping Between Models}
\label{sec:material-history-mapping}
%======================================================================

The most subtle and potentially impactful aspect of the multiscale
state transfer concerns the \emph{material internal variables}
$\boldsymbol{\alpha}$~that encode loading history.  Currently, each
scale maintains its own independent material history:

\begin{itemize}
  \item The global beam model uses fiber-section integration with
    Kent--Park concrete and Menegotto--Pinto steel.  Each fiber
    accumulates its own internal state
    $\boldsymbol{\alpha}^{\text{macro}}$.
  \item The sub-model uses 3D Ko--Bathe concrete hexahedra with
    embedded Menegotto--Pinto truss elements for rebar.  Each Gauss
    point / truss element maintains its own
    $\boldsymbol{\alpha}^{\text{micro}}$.
\end{itemize}

When the sub-model is first activated at a given time step, the
micro-scale material points start from a virgin state, regardless of
the loading history that the global model has already undergone.
This inconsistency leads to two problems:

\begin{enumerate}
  \item \textbf{Initial transient:}  The sub-model must ``catch up''
    to the global solution, effectively re-experiencing the loading
    history incrementally.  This slows convergence and may introduce
    spurious damage.
  \item \textbf{Incorrect effective tangent:}  The sub-model's
    tangent stiffness $\mathbb{D}_{\text{hom}}$ starts from the
    virgin elastic state rather than from the degraded/hardened state
    of the macro-scale, causing a stiffness mismatch that impairs
    the global Newton convergence.
\end{enumerate}

% ─────────────────────────────────────────────────────────────────────
\subsection{Menegotto--Pinto History Transfer: Fibers to Trusses}
\label{ssec:menegotto-transfer}

The Menegotto--Pinto steel model maintains the following internal
state at each integration point:

\begin{equation}
  \boldsymbol{\alpha}_{\text{MP}} = \bigl\{
    \varepsilon_r,\; \sigma_r,\;
    \varepsilon_0,\; \sigma_0,\;
    \varepsilon_r^{\text{prev}},\; \sigma_r^{\text{prev}},\;
    \varepsilon_{\max},\; \xi,\;
    d,\; \varepsilon_c,\; \sigma_c,\;
    \text{yielded}
  \bigr\},
  \label{eq:mp-state}
\end{equation}

where $(\varepsilon_r, \sigma_r)$ is the last reversal point,
$(\varepsilon_0, \sigma_0)$ the asymptote intersection,
$\xi = |\varepsilon_0 - \varepsilon_r|$ the shift parameter for the
R-evolution $R(\xi) = R_0 - c_{R1}\,\xi/(c_{R2} + \xi)$,
$d \in \{+1, -1, 0\}$ the loading direction,
and $(\varepsilon_c, \sigma_c)$ the last committed strain and stress.

In the global beam model, this state lives at each \emph{fiber}
integration point.  In the sub-model, the same Menegotto--Pinto
constitutive law governs each \emph{truss element} representing
a reinforcing bar.  The key observation is that:

\begin{quote}
\itshape
A macro-scale fiber at position $(y_i, z_i)$ in the section, carrying
area $A_i$, has a one-to-one correspondence with a micro-scale truss
element at the same geometric position, with the same cross-sectional
area and the same material parameters.
\end{quote}

Therefore, the history transfer is, in principle, a direct copy:
\begin{equation}
  \boldsymbol{\alpha}_{\text{MP}}^{\text{truss}_j}
    \leftarrow
  \boldsymbol{\alpha}_{\text{MP}}^{\text{fiber}_i},
  \qquad
  \text{where } j = \mathcal{M}(i),
  \label{eq:mp-transfer}
\end{equation}
where $\mathcal{M}: \{\text{fibers}\} \to \{\text{trusses}\}$ is the
geometric mapping that associates each fiber with its corresponding
truss element (by nearest-neighbour matching of the $(y,z)$ coordinates
in the section plane).

For the material parameters $(E_0, f_y, b, R_0, c_{R1}, c_{R2})$
to be identical at both scales (which they are in the current
implementation), the transfer is exact: the truss element, initialized
with the fiber's reversal points and accumulated plastic strain, will
produce the identical stress--strain response for any subsequent
loading path.

% ─────────────────────────────────────────────────────────────────────
\subsection{Concrete History Transfer: Fibers to Hexahedra}
\label{ssec:concrete-transfer}

The concrete case is more complex because the constitutive models at
the two scales differ:

\begin{itemize}
  \item \textbf{Macro-scale:}  Kent--Park uniaxial model (confined /
    unconfined), parameterised by $f'_c$, $\varepsilon_0$,
    $\varepsilon_{20}$, and damage scalars.
  \item \textbf{Micro-scale:}  Ko--Bathe 3D multiaxial model with crack
    tensors per Gauss point, tracking crack opening, closing, and
    re-opening in up to three directions.
\end{itemize}

A direct variable-by-variable copy is not possible.  Instead, the
transfer must be \emph{semantic}:

\begin{enumerate}
  \item \textbf{Damage scalar mapping:}  The Kent--Park damage
    parameter $d^{\text{macro}}(\varepsilon)$ can be mapped to an
    equivalent isotropic damage in the Ko--Bathe model by matching
    the secant stiffness degradation:
    \begin{equation}
      (1 - d_{\text{KB}})\,E_0^{\text{3D}}
        \approx (1 - d_{\text{KP}})\,E_0^{\text{1D}}.
      \label{eq:damage-matching}
    \end{equation}
  \item \textbf{Crack direction:}  The 1D fiber model does not track
    crack orientation.  A reasonable initialisation is to set the
    principal crack direction perpendicular to the fiber's axial
    direction (i.e.\ normal to the beam axis), consistent with
    the flexural crack pattern expected at that section location.
  \item \textbf{Strain history:}  The maximum compressive and tensile
    strains experienced by each fiber can be used to initialise the
    Ko--Bathe equivalent strain thresholds, ensuring that the
    micro-scale model starts with the correct damage level.
\end{enumerate}

% ─────────────────────────────────────────────────────────────────────
\subsection{Proposed Architecture: StateTransferProtocol}
\label{ssec:state-transfer-protocol}

To formalize and extensibly implement the above transfers, we propose
a C++20 concept:

\begin{lstlisting}
template <typename T>
concept StateTransferProtocol = requires(
    T& transfer,
    const GlobalModel& macro,
    LocalModel& micro,
    std::size_t gauss_idx)
{
    // Map fiber internal state to truss element
    transfer.map_rebar_history(macro, micro);

    // Inject reconstructed stress/strain at Gauss points
    transfer.inject_initial_stress(macro, micro);

    // Inject reconstructed strain at Gauss points
    transfer.inject_initial_strain(macro, micro);

    // Map concrete damage from 1D fibers to 3D Gauss points
    transfer.map_concrete_damage(macro, micro);
};
\end{lstlisting}

A concrete implementation would proceed in four stages:

\begin{enumerate}
  \item \textbf{Geometric mapping:}  Establish the correspondences
    fiber~$\leftrightarrow$~truss and fiber~$\leftrightarrow$~Gauss
    point using a spatial search in the section plane.
  \item \textbf{Rebar history injection:}  Copy
    $\boldsymbol{\alpha}_{\text{MP}}^{\text{fiber}}$ verbatim into the
    matched truss's \texttt{MenegottoPintoState}.
  \item \textbf{Concrete state seeding:}  Evaluate the reconstructed
    strain field at each Gauss point and initialise the Ko--Bathe
    crack tensor from the macro-scale damage.
  \item \textbf{Stress initialisation:}  Set
    $\boldsymbol{\sigma}^{(0)}_{\text{GP}}
      = \boldsymbol{\sigma}^{\text{macro}}(y_{\text{GP}}, z_{\text{GP}})$
    and use it as the initial stress for the sub-model residual.
\end{enumerate}


%======================================================================
\section{Expected Impact on Convergence}
\label{sec:convergence-impact}
%======================================================================

The completeness of the state transfer directly affects the convergence
behaviour of the sub-model nonlinear solve.  We can distinguish three
levels of transfer and their expected effects:

\begin{table}[htbp]
\centering
\caption{State transfer levels and expected convergence impact.}
\label{tab:transfer-levels}
\begin{tabular}{@{}clp{5.5cm}p{4.5cm}@{}}
\toprule
\textbf{Level} & \textbf{Transfer Content} & \textbf{Effect on Convergence} & \textbf{Status} \\
\midrule
0 & Displacement BCs only
  & Sub-model starts from virgin material; Newton must traverse the
    entire loading path.  High iteration count.
  & Implemented \\
1 & + Reconstructed $\boldsymbol{\varepsilon}$, $\boldsymbol{\sigma}$
  & Initial stress field reduces residual norm at first Newton
    iteration.  Fewer iterations.
  & Implemented (\S\ref{sec:state-injection-impl}). \\
2 & + Rebar history ($\boldsymbol{\alpha}_{\text{MP}}$)
  & Truss elements start on the correct loading branch with the
    correct tangent modulus.
    Eliminates spurious elastic unloading of reinforcement.
  & Implemented (\S\ref{sec:state-injection-impl}). \\
3 & + Concrete damage mapping
  & Ko--Bathe Gauss points start with degraded stiffness consistent
    with macro-scale damage.  Prevents artificial stiffness
    recovery at the transition.
  & Proposed; requires semantic damage mapping. \\
\bottomrule
\end{tabular}
\end{table}

The hypothesis is that the progression from Level~0 to Level~3
will reduce the required number of adaptive sub-steps and bisection
depth in the sub-model solver, with the most significant improvement
coming from Level~2 (rebar history), since the steel constitutive
response is highly path-dependent and the tangent modulus undergoes
order-of-magnitude changes between loading and unloading branches.

The physical intuition is straightforward:
\begin{quote}
\itshape
A sub-model that starts from a physical state consistent with the
global solution does not need to ``rediscover'' the loading history.
It begins at the right point on the right loading branch, with the
right stiffness, and needs only to accommodate the \emph{incremental}
change from the previous global step---not the entire accumulated
deformation.
\end{quote}

This is precisely analogous to the benefit of using the converged
solution as the initial guess for Newton's method: starting close to
the answer means fewer iterations to convergence.

%======================================================================
\section{Implementation Roadmap}
\label{sec:impl-roadmap}
%======================================================================

The following concrete steps are proposed (items marked~\checkmark have
been completed; see \S\ref{sec:state-injection-impl} for details):

\begin{enumerate}
  \item[\checkmark] \textbf{Add \texttt{set\_internal\_state()} to
    constitutive relations} and the \texttt{ExternallyStateInjectable}
    concept.  Extend the \texttt{Material<>} vtable with
    \texttt{StateRef}-based injection using
    \texttt{WrapsInjectableRelation} dispatch.
  \item[\checkmark] \textbf{Implement fiber$\,\leftrightarrow\,$truss
    geometric mapping:}  \texttt{FiberToTrussMapper.hh} with
    nearest-neighbour matching and batch state injection.
  \item[\checkmark] \textbf{Implement
    \texttt{reconstruct\_3d\_strain()} and
    \texttt{reconstruct\_3d\_stress()}:}  Poisson-enriched strain
    field and nonlinear fiber-based stress reconstruction in
    \texttt{FieldTransfer.hh}.
  \item[\checkmark] \textbf{Hook in sub-model evolver:}
    \texttt{StateTransferCallback} invoked in \texttt{first\_solve()}
    before the incremental ramp.
  \item \textbf{Implement \texttt{StateTransferProtocol}:}  A
    concrete orchestrator class that wires all components together
    and provides the four-stage transfer described in
    \S\ref{ssec:state-transfer-protocol}.
  \item \textbf{Benchmark:}  Compare Newton iteration counts and
    sub-step requirements for the cyclic validation column at
    all four transfer levels, using the Full Extended~(50\,mm)
    protocol.
\end{enumerate}

The critical implementation constraint is that the state injection
must respect the \texttt{IncrementalControlPolicy} idempotency
invariant: the injected state is set once at the beginning of the
sub-model solve and does not depend on the bisection parameter~$p$.
It is an \emph{initial condition}, not a boundary condition.


%======================================================================
\section{Implementation of State Injection Infrastructure}
\label{sec:state-injection-impl}
%======================================================================

This section documents the concrete implementation of Levels~1 and~2
from Table~\ref{tab:transfer-levels}, covering the material state
injection API, the fiber-to-truss geometric mapping, and the
stress/strain reconstruction pipeline.  The implementation was
completed in three phases (A, B, C) as detailed below.


% ─────────────────────────────────────────────────────────────────────
\subsection{Phase A: Material State Injection API}
\label{ssec:phase-a-injection-api}
% ─────────────────────────────────────────────────────────────────────

The foundation of the state transfer mechanism is a new concept and
a virtual dispatch chain that enables injection of material internal
variables through the type-erasure boundary.

\subsubsection{The \texttt{ExternallyStateInjectable} Concept}

A new Level~4 requirement extends the concept hierarchy in
\texttt{ConstitutiveRelation.hh}:

\begin{lstlisting}
template <typename R>
concept ExternallyStateInjectable =
    InelasticConstitutiveRelation<R> &&
    requires(R r, const typename R::InternalVariablesT& alpha) {
        { r.set_internal_state(alpha) };
    };
\end{lstlisting}

Both \texttt{MenegottoPintoSteel} and \texttt{KoBatheConcrete3D}
satisfy this concept via their new \texttt{set\_internal\_state()}
methods, which replace the internal state~$\boldsymbol{\alpha}$
with an externally provided snapshot:
\[
  \boldsymbol{\alpha} \leftarrow \boldsymbol{\alpha}_{\text{ext}}.
\]

\subsubsection{The \texttt{StateRef} Non-Owning Type-Erased Reference}

A key design decision is how to pass the typed internal state through
the type-erased \texttt{Material<>} wrapper.  The na\"ive approach of
using \texttt{std::any} incurs a heap allocation per injection
(the state structs---particularly \texttt{KoBatheState3D} with its
Eigen vectors---exceed the small-buffer optimisation threshold) plus
an RTTI comparison on extraction via \texttt{std::any\_cast}.

Since the injected state is always an lvalue whose lifetime exceeds
the injection call, we use a zero-copy alternative:

\begin{lstlisting}
struct StateRef {
   const void*           data_  = nullptr;
   const std::type_info* type_  = nullptr;

   template <typename T>
   static StateRef from(const T& value) noexcept {
      return StateRef{&value, &typeid(T)};
   }

   template <typename T>
   const T& as() const {
      if (!type_ || *type_ != typeid(T))
         throw std::runtime_error("StateRef type mismatch");
      return *static_cast<const T*>(data_);
   }
};
\end{lstlisting}

\texttt{StateRef} is a non-owning reference that stores only a
\texttt{const void*} pointer and a \texttt{const std::type\_info*}
for runtime type safety.  Compared to \texttt{std::any}:

\begin{itemize}
  \item \textbf{Zero allocation:}  No heap allocation regardless of
    state size.
  \item \textbf{Zero copy:}  The state struct is never copied;
    \texttt{StateRef::from()} stores a pointer to the original.
  \item \textbf{Type-safe:}  \texttt{as<T>()} performs an RTTI
    comparison (\texttt{typeid} equality) and throws on mismatch,
    providing the same safety guarantee as \texttt{std::any\_cast}.
  \item \textbf{Trivial overhead:}  Two pointers (16~bytes on 64-bit)
    passed by value through the virtual call chain.
\end{itemize}


\subsubsection{Virtual Dispatch Through the Type-Erasure Chain}

The injection traverses four layers of the type-erasure architecture:

\begin{center}
\small
\begin{tabular}{@{}lll@{}}
\toprule
\textbf{Layer} & \textbf{Class} & \textbf{Action} \\
\midrule
1 & \texttt{Material<P>} / \texttt{MaterialRef<P>}
  & Forwards to \texttt{pimpl\_->inject\_internal\_state(state)} \\
2 & \texttt{MaterialConcept} (virtual base)
  & Dispatches to concrete model override \\
3 & \texttt{OwningMaterialModel<MT, Int>}
  & \texttt{if constexpr} branch at compile time \\
4 & Concrete material / \texttt{ConstitutiveSite}
  & Calls \texttt{set\_internal\_state()} \\
\bottomrule
\end{tabular}
\end{center}

At Layer~3, two concepts control the dispatch:

\begin{itemize}
  \item \texttt{HasSetInternalState<MT>}: the stored type directly
    provides \texttt{set\_internal\_state()}.  Applies when the
    \texttt{Material<>} wraps a raw constitutive relation.
  \item \texttt{WrapsInjectableRelation<MT>}: the stored type is a
    \texttt{ConstitutiveSite} that exposes \texttt{constitutive\_relation()}
    returning an injectable relation.  Applies when the \texttt{Material<>}
    wraps an \texttt{InelasticMaterial<R>}.
\end{itemize}

For the \texttt{WrapsInjectableRelation} path, the injection must
update \emph{both} the raw relation's internal state and the
\texttt{ConstitutiveSite}'s \texttt{algorithmic\_state()}, which is
a separate copy used by the Level~3 (externally-driven) dispatch:
\begin{lstlisting}
// WrapsInjectableRelation path:
material_.constitutive_relation().set_internal_state(s);
material_.algorithmic_state() = s;  // sync Level 3 dispatch state
\end{lstlisting}

This dual-sync is necessary because
\texttt{MaterialInstance::compute\_response(k)} delegates to
\texttt{relation\_->compute\_response(k, algorithmic\_state())},
not to the relation's own embedded state.


\subsubsection{Element-Level Injection Methods}

\texttt{TrussElement} and \texttt{ContinuumElement} expose
element-level injection methods that forward through
\texttt{MaterialPoint} to the type-erased material:

\begin{lstlisting}
// TrussElement: single material point
void inject_material_state(impl::StateRef state);

// ContinuumElement: per Gauss point
void inject_material_state(std::size_t gp, impl::StateRef state);
void inject_material_state(const std::vector<impl::StateRef>& states);
\end{lstlisting}


% ─────────────────────────────────────────────────────────────────────
\subsection{Phase B: Fiber-to-Truss Geometric Mapping}
\label{ssec:phase-b-fiber-truss}
% ─────────────────────────────────────────────────────────────────────

The \texttt{FiberSection} in the macro-beam stores per-fiber
\texttt{Fiber\{y, z, A, Material<UniaxialMaterial>\}} records.  The
micro-model's reinforcement is a set of \texttt{TrussElement<3>}
placed at matching $(y, z)$ positions.  The mapper establishes a
bijection:

\[
  \mathcal{M}: \{\text{steel fibers}\} \to \{\text{truss bars}\},
  \quad
  \mathcal{M}(i) = \arg\min_{j} \|
    (y^{\text{fiber}}_i, z^{\text{fiber}}_i) -
    (y^{\text{bar}}_j, z^{\text{bar}}_j)
  \|_2,
  \quad
  \text{s.t.}\ \|\cdot\| < \texttt{tol}.
\]

\texttt{build\_fiber\_truss\_map()} identifies steel fibers (those
satisfying \texttt{supports\_state\_injection()}), performs
nearest-neighbour matching, and produces a vector of
\texttt{FiberTrussLink\{fiber\_index, bar\_index, distance\}} records.

The injection function \texttt{inject\_rebar\_states()} takes a vector
of \texttt{RebarStatePacket\{bar\_index, MenegottoPintoState\}} and
injects each state into the corresponding truss element via
\texttt{StateRef::from()}.


% ─────────────────────────────────────────────────────────────────────
\subsection{Phase C: Stress/Strain Reconstruction and Hook}
\label{ssec:phase-c-reconstruction}
% ─────────────────────────────────────────────────────────────────────

Two reconstruction functions in \texttt{FieldTransfer.hh} extend the
existing beam-to-continuum field transfer:

\begin{enumerate}
  \item \textbf{\texttt{reconstruct\_3d\_strain()}:}  Enriches the
    beam kinematic field with Poisson effects:
    \[
      \varepsilon_{yy} = \varepsilon_{zz} = -\nu\,\varepsilon_{xx},
    \]
    producing a full 6-component Voigt vector at each $(y,z)$ position.

  \item \textbf{\texttt{reconstruct\_3d\_stress()}:}  Finds the
    nearest fiber in the \texttt{FiberSection}, queries the fiber's
    nonlinear uniaxial stress via \texttt{compute\_response()}, and
    assembles a full stress tensor with elastic shear fallback:
    $\sigma_{xy} = G\gamma_y$, $\sigma_{xz} = G\gamma_z$.
\end{enumerate}

Both functions operate on the converged macro-beam state and do not
require iteration.

\subsubsection{Integration Hook in the Sub-Model Evolver}

The \texttt{NonlinearSubModelEvolver} gains a
\texttt{StateTransferCallback} member:

\begin{lstlisting}
using StateTransferCallback =
    std::function<void(typename MixedModel::container_type&)>;

void set_state_transfer(StateTransferCallback cb);
\end{lstlisting}

The callback is invoked in \texttt{first\_solve()}, between the SNES
setup (Step~3) and the incremental ramp (Step~4).  At this point the
sub-model mesh is fully assembled but the Newton solve has not yet
begun, making it the correct location for one-time state initialisation
that respects the idempotency invariant.


% ─────────────────────────────────────────────────────────────────────
\subsection{Verification}
\label{ssec:state-injection-verification}
% ─────────────────────────────────────────────────────────────────────

The implementation is verified by a dedicated unit test
(\texttt{test\_state\_injection.cpp}) that exercises four scenarios:

\begin{enumerate}
  \item \textbf{MenegottoPintoSteel roundtrip:}  Cycle the material
    through tension and compression reversals, extract
    $\boldsymbol{\alpha}$, inject into a fresh instance, verify
    identical stress and tangent response.
  \item \textbf{KoBatheConcrete3D roundtrip:}  Same pattern for the
    3D concrete model under uniaxial compression.
  \item \textbf{Type-erased injection:}  Create a
    \texttt{Material<UniaxialMaterial>} wrapping an
    \texttt{InelasticMaterial<MenegottoPintoSteel>} (the full
    \texttt{ConstitutiveSite} path), inject via
    \texttt{StateRef}, verify the response matches the raw
    relation.
  \item \textbf{Concept satisfaction:}  Compile-time
    \texttt{static\_assert} for the \texttt{ExternallyStateInjectable}
    concept on both materials.
\end{enumerate}

All four tests pass.  The full existing test suite (71~tests) also
passes with zero regressions.


% ─────────────────────────────────────────────────────────────────────
\subsection{Summary of Modified Files}
\label{ssec:modified-files}
% ─────────────────────────────────────────────────────────────────────

Table~\ref{tab:modified-files} lists every file modified or created
for the state injection infrastructure.

\begin{table}[htbp]
\centering
\caption{Files modified or created for the state injection infrastructure.}
\label{tab:modified-files}
\small
\begin{tabular}{@{}lp{7.5cm}@{}}
\toprule
\textbf{File} & \textbf{Change} \\
\midrule
\texttt{ConstitutiveRelation.hh}
  & Added \texttt{ExternallyStateInjectable} concept \\
\texttt{MenegottoPintoSteel.hh}
  & Added \texttt{set\_internal\_state()} \\
\texttt{KoBatheConcrete3D.hh}
  & Added \texttt{set\_internal\_state()} \\
\texttt{Material.hh}
  & Added \texttt{StateRef}, \texttt{HasSetInternalState},
    \texttt{WrapsInjectableRelation} concepts; extended vtable
    with \texttt{inject\_internal\_state(StateRef)} \\
\texttt{MaterialPoint.hh}
  & Forwarding to \texttt{material\_.inject\_internal\_state()} \\
\texttt{TrussElement.hh}
  & Added \texttt{inject\_material\_state(StateRef)} \\
\texttt{ContinuumElement.hh}
  & Added \texttt{inject\_material\_state()} (single GP + batch) \\
\texttt{FiberToTrussMapper.hh} \emph{(new)}
  & Nearest-neighbour fiber$\leftrightarrow$truss matching;
    \texttt{RebarStatePacket}; \texttt{inject\_rebar\_states()} \\
\texttt{FieldTransfer.hh}
  & Added \texttt{reconstruct\_3d\_strain()},
    \texttt{reconstruct\_3d\_stress()} \\
\texttt{NonlinearSubModelEvolver.hh}
  & Added \texttt{StateTransferCallback} hook in
    \texttt{first\_solve()} \\
\texttt{test\_state\_injection.cpp} \emph{(new)}
  & Unit test: roundtrip for both materials + type-erased path \\
\bottomrule
\end{tabular}
\end{table}
